\topic{本章实验、故意破坏与练习}

\begin{experimentbox}{五个 FLUSH 点分别重启}
在 data、prepared、commit、checkpoint、clear 五个 FLUSH 后复制磁盘并重启，
记录 journal 字段、replay 次数、路径结果和 fsck。
\end{experimentbox}

\begin{faultinjectionbox}{commit 先于 prepared 稳定}
交换 prepared FLUSH 与 commit。预期可能留下 valid commit 和不完整 payload。
恢复为 payload 全稳定后才写 commit。
\end{faultinjectionbox}

\begin{pitfallbox}
IRQ14 只表示请求完成；FLUSH 决定持久顺序；checkpoint 失败时不能清 journal。
\end{pitfallbox}

\textbf{练习}
\begin{enumerate}[leftmargin=2.2em]
  \item commit 缺失但 payload 完整时选择哪种状态？
  \item checkpoint 一半后断电怎样恢复？
  \item 同一 target stage 两次消耗几个 credit？
\end{enumerate}
\begin{answerbox}
1. 旧状态；2. 重放全部 target；3. 一个。
\end{answerbox}
